% Définitions des résumés (projet DevSecOps PaaS — microservices / plate-forme Kubernetes)
% Inclus depuis rapport.tex avant \begin{document}

\newcommand{\frenchabstract}{%
Ce rapport documente la plate-forme logicielle développée dans le dépôt \texttt{paas/} pour notre projet de fin d'études : un~\emph{plan de contrôle} unique en Next.js~(App Router) qui regroupe une interface métier et des routes API~(Prisma\,/\,PostgreSQL, sessions JWT) pour créer et suivre les projets, déclencher des builds~(via \texttt{BUILD\_BACKEND}=\texttt{jenkins} ou~\texttt{tekton}), suivre promotions GitOps~\texttt{GITOPS\_*} et statuts jusqu'à Argo~CD et au cluster Kubernetes.

Le moteur \emph{BuildPlanner} associe chaque dépôt à un profil de build géré~(\texttt{node}, \texttt{python}, \texttt{java}, \texttt{static}, \texttt{custom}) et choisit modèle de plate-forme ou contrat Dockerfile. Les exécutions Jenkins ou les \texttt{PipelineRun} Tekton alimentent Harbor~; la promotion pousse l'image~(et le digest optionnel) vers le dépôt GitOps. Le flux et la documentation \texttt{paas/README.md} et \texttt{paas/docs/} alignent ce mémoire sur le code réel.

La stack sécurité et supervision~(SonarQube, Dependency-Track, Trivy, Cosign, politiques Kyverno~/~OPA selon configuration) se branche sur les URL et jetons lus par le serveur : le périmètre est celui de l'application que nous construisons, pas un catalogue théorique.
}

\newcommand{\frenchkeywords}{%
DevSecOps, plate-forme as a Service, Kubernetes, GitOps, Jenkins, Tekton, Argo CD, Harbor, sécurité des conteneurs, CI/CD%
}

\newcommand{\englishabstract}{%
This thesis documents the software we built in the \texttt{paas/} repository: a single Next.js App Router \emph{control plane} with Prisma/PostgreSQL persistence, JWT sessions, and API routes that onboard projects, trigger builds through a pluggable backend (\texttt{BUILD\_BACKEND} is \texttt{jenkins} or \texttt{tekton}), push image references to a GitOps repo (\texttt{GITOPS\_*}), and track deployment status through Argo CD and Kubernetes.

The \emph{BuildPlanner} maps repositories to managed profiles (\texttt{node}, \texttt{python}, \texttt{java}, \texttt{static}, \texttt{custom}) and chooses platform templates or a custom Dockerfile contract. Jenkins jobs or Tekton \texttt{PipelineRun} resources feed Harbor; promotion is artifact-first. The narrative matches \texttt{paas/README.md}, \texttt{paas/docs/STACK\_OVERVIEW.md}, and \texttt{paas/docs/DEVSOPS\_PAAS\_ARCHITECTURE.md}.

Security and observability tools (SonarQube, Dependency-Track, Trivy, Cosign, policy engines as configured via environment variables) are integrated as adapters in our application—not a generic toolchain survey divorced from implementation.
}

\newcommand{\englishkeywords}{%
DevSecOps, Platform as a Service, Kubernetes, GitOps, Jenkins, Tekton, Argo CD, Harbor, container security, CI/CD%
}

\newcommand{\arabicabstract}{%
يقدّم هذا التقرير تصميم منصّة برمجيات تهدف إلى دمج الأمن ضمن سلسلة تطوير ونشر مستمرّة مستهدفاً بيئة تجمّع لتشغيل تطبيقات حاوويّة. تتيح المنصّة تسجيل المشاريع، وتخطيط البناء وفق ملفّات مهيّأة أو قوالب نظامية، وتخزين صور الحاويات في سجل مركزي، ثم ضبط الحالة المرغوبة عبر آلية شبيهة بإدارة التهيئة من المستودع، ومتابعة نشر الطبقات وفق المراحل الموحدة. كما تجمع بين أدوات تحليل جودة الكود وحصر المكوّنات وفحص الثغرات والسياسات لدى الاستقبال وفق مقاربة «تطوير آمن للعمليات».%
}

\newcommand{\arabickeywords}{%
تطوير آمن للعمليات، تجمّع خوادم، إدارة التهيئة من المستودع، خط أنابيب البرمجيات، الحاويات، الأمن%
}
